\section{Shortest Distance in a Binary Maze}
\label{sec:graph-shortest-binary-maze}

\problemstatement{Given an $n\times m$ grid of $0$s (blocked) and $1$s
(open), a source cell and a destination cell, find the length of the
shortest path between them moving only through open cells in the four
cardinal directions. Return $-1$ if unreachable.}

\begin{patternbox}
A grid where every move costs the same one step is a
\textbf{unit-weight} graph, so plain \textbf{BFS} finds the shortest
path -- exactly the case from
Section~\ref{sec:graph-shortest-path-unit-weights}, just with the
graph given implicitly as a grid rather than an adjacency list. No
Dijkstra's is needed when all moves cost $1$.
\end{patternbox}

\subsection*{The grid is a graph in disguise}

Each open cell is a node; each pair of side-adjacent open cells is an
edge of weight $1$. Because BFS explores strictly in order of increasing
number of steps, the first time it reaches the destination it has done
so by a shortest path. We store $(\textit{row},\textit{col},
\textit{distance})$ in the queue (or track distance in a separate grid)
and expand the four neighbors, skipping blocked or out-of-bounds cells
and cells already visited.

\begin{ideabox}[BFS Layers Are Distance Layers]
Since every edge weighs $1$, BFS's $k$-th layer holds exactly the cells
at distance $k$ from the source. Marking a cell visited the moment it is
enqueued guarantees each cell is reached by its shortest distance and
never reprocessed.
\end{ideabox}

\subsection*{Algorithm}

\begin{enumerate}
  \item If the source is blocked, return $-1$ (or $0$ if source equals
        destination). Push the source with distance $0$; mark it
        visited.
  \item While the queue is non-empty, pop $(r,c,d)$; if it is the
        destination, return $d$.
  \item For each of the four neighbors that is in bounds, open, and
        unvisited: mark visited and push it with distance $d+1$.
  \item If the queue empties without reaching the destination, return
        $-1$.
\end{enumerate}

\subsection*{C++ implementation}

\begin{lstlisting}
#include <bits/stdc++.h>
using namespace std;

int shortestPathBinaryMaze(vector<vector<int>>& grid,
                           pair<int,int> src, pair<int,int> dest) {
    int n = grid.size(), m = grid[0].size();
    if (grid[src.first][src.second] == 0) return -1;
    if (src == dest) return 0;

    vector<vector<bool>> vis(n, vector<bool>(m, false));
    queue<tuple<int,int,int>> q;             // {row, col, distance}
    q.push({src.first, src.second, 0});
    vis[src.first][src.second] = true;
    int dr[] = {-1, 1, 0, 0}, dc[] = {0, 0, -1, 1};

    while (!q.empty()) {
        auto [r, c, d] = q.front(); q.pop();
        for (int k = 0; k < 4; k++) {
            int nr = r + dr[k], nc = c + dc[k];
            if (nr < 0 || nr >= n || nc < 0 || nc >= m) continue;
            if (grid[nr][nc] == 0 || vis[nr][nc]) continue;
            if (nr == dest.first && nc == dest.second) return d + 1;
            vis[nr][nc] = true;
            q.push({nr, nc, d + 1});
        }
    }
    return -1;
}

int main() {
    vector<vector<int>> grid = {
        {1, 1, 1, 1},
        {1, 0, 0, 1},
        {1, 1, 1, 1}
    };
    cout << shortestPathBinaryMaze(grid, {0,0}, {2,3}) << endl;  // 5
    return 0;
}
\end{lstlisting}

\subsection*{Dry run}

\dryrun{Source $(0,0)$, destination $(2,3)$ on the $3\times4$ grid
above.}

\begin{itemize}
  \item Layer $0$: $(0,0)$. Layer $1$: $(1,0),(0,1)$.
  \item Layer $2$: $(2,0),(0,2)$. Layer $3$: $(2,1),(0,3)$.
  \item Layer $4$: $(2,2),(1,3)$. Layer $5$: expanding $(2,2)$ or
        $(1,3)$ reaches $(2,3)$ -- return $5$.
  \item The blocked cells $(1,1),(1,2)$ force the path around the
        edges.
\end{itemize}

\subsection*{Complexity}

\begin{complexitybox}
\textbf{Time Complexity.} $O(n\cdot m)$ -- each cell is enqueued and
processed at most once. \textbf{Space Complexity.} $O(n\cdot m)$ for the
visited grid and the BFS queue.
\end{complexitybox}

\begin{takeawaybox}
When every move on a grid costs the same, BFS -- not Dijkstra's -- is the
right tool, and the grid is simply an implicit unit-weight graph. Save
Dijkstra's for grids where cells carry differing move costs, such as
Path With Minimum Effort (Section~\ref{sec:graph-path-min-effort}).
\end{takeawaybox}
